\documentclass[12pt]{article}                                                                                                           
\usepackage[utf8]{inputenc}                                                                                                             
\usepackage[T1]{fontenc}                                                                                                                
\usepackage{amsmath, amssymb}                                                                                                           
\usepackage{geometry}                                                                                                                   
\usepackage{xcolor}                                                                                                                     
\usepackage{lipsum}                                                                                                                     
\usepackage{titlesec}                                                                                                                   
\usepackage{fancyhdr}                                                                                                                   
\usepackage[most]{tcolorbox}                                                                                                                  
\usepackage{graphicx}                                                                                                                   
\usepackage{float}                                                                                      
\usepackage{tikz}
\usetikzlibrary{shapes.geometric, arrows.meta}
\usetikzlibrary{arrows.meta, decorations.markings}                                                                                      
                                                                                                                                        
\geometry{a4paper, margin=2.5cm}                                                                                                        
\pagestyle{fancy}                                                                                                                       
\fancyhf{}                                                                                                                              
\rhead{Problem Set II}                                                                                                                   
\lhead{Probability and Statistics}                                                                                                                     
\cfoot{\thepage}                                                                                                                        
                                                                                                                                        
% Fancy titles                                                                                                                          
\titleformat{\section}{\normalfont\Large\bfseries}{Exercise \thesection:}{1em}{}                                                        
\titleformat{\subsection}{\normalfont\bfseries}{Correction:}{1em}{}                                                                     
                                                                                                                                        
% Custom box for correction with breakable option
\newtcolorbox{correctionbox}{                                                                                                           
  colback=gray!5,                                                                                                                       
  colframe=black,                                                                                                                       
  fonttitle=\bfseries,                                                                                                                  
  title=Correction,                                                                                                                     
  before skip=10pt,                                                                                                                     
  after skip=10pt,
  breakable,
  enhanced
}                                                                                                                                       
                                                                                                                                        
% Fix headheight warning
\setlength{\headheight}{14.49998pt}

\begin{document}

\begin{center}
\Large\textbf{Ibn Tofail University} \\[1em]
\large\textit{Probability and Statistics} \\[2em]
\large\textit{Problem Set II} \\[0.5em]
\end{center}

\vspace{1cm}

% ----------- EXERCISE 1 -------------
\section{}
Consider a cubic die whose faces are numbered from 1 to 6, biased such that the probability of obtaining face number \( k \) is proportional to \( k \), with proportionality coefficient \( p \). We roll the die once and denote by \( X \) the number on the face obtained.
\begin{enumerate}
    \item Determine the probability distribution of \( X \) (as a function of \( p \)).
    \item Determine the value of \( p \).
    \item Calculate \( E(X) \).
    \item Let \( Y = \frac{1}{X} \). Determine the distribution of \( Y \) and \( E(Y) \).
\end{enumerate}

\begin{correctionbox}

\end{correctionbox}

% ----------- EXERCISE 2 -------------
\section{}
A straight track is divided into squares numbered 0, 1, 2, ..., \( n \), ... from left to right. A flea moves to the right, randomly jumping 1 or 2 squares at each move. Initially, it is on square 0. Let \( X_n \) be the random variable equal to the number of the square occupied by the flea after \( n \) jumps.
\begin{enumerate}
    \item Determine the probability distribution of \( X_1 \), and calculate \( E(X_1) \) and \( V(X_1) \).
    \item Let \( Y_n \) be the random variable equal to the number of times the flea jumped one square during the first \( n \) jumps. Determine the distribution of \( Y_n \), then \( E(Y_n) \) and \( V(Y_n) \).
    \item Express \( X_n \) as a function of \( Y_n \) and deduce the probability distribution of \( X_n \), then \( E(X_n) \) and \( V(X_n) \).
\end{enumerate}

\begin{correctionbox}

\end{correctionbox}

% ----------- EXERCISE 3 -------------
\section{}
We play heads or tails with a biased coin. In each toss, the probability of getting tails is \( \frac{1}{3} \). The tosses are assumed to be independent. Let \( X \) be the random variable equal to the number of tosses needed to obtain two consecutive heads for the first time. Let \( n \) be a non-zero natural integer. We denote \( p_n \) as the probability of the event \( [X = n] \). Additionally, we denote \( F_i \) as the event "obtaining tails on the \( i \)-th toss".
\begin{enumerate}
    \item Express the events \( [X = 2] \), \( [X = 3] \), \( [X = 4] \), \( [X = 5] \) in terms of the events \( F_i \) and \( \overline{F_i} \). Determine the values of \( p_1, p_2, p_3, p_4, p_5 \).
    \item Using the total probability formula, distinguishing two cases based on the result of the first toss, show that:
    \[
    \forall n \geq 3, \quad p_n = \frac{2}{9} p_{n-2} + \frac{1}{3} p_{n-1}
    \]
    \item Deduce the expression of \( p_n \) as a function of \( n \), for \( n \geq 1 \).
    \item Calculate \( E(X) \).
\end{enumerate}

\begin{correctionbox}
    
\end{correctionbox}

% ----------- EXERCISE 4 -------------
\section{}
An urn contains 2 white balls and \( n - 2 \) red balls. We perform \( n \) draws without replacement from this urn. Let \( X \) be the rank (draw number) of the first white ball drawn and \( Z \) be the rank of the second white ball drawn.
\begin{enumerate}
    \item Determine the joint distribution of the pair \( (X, Z) \).
    \item Deduce the distribution of \( Z \).
\end{enumerate}

\begin{correctionbox}

\end{correctionbox}


% ----------- EXERCISE 5 -------------
\section{}
We have \( n \) boxes numbered from 1 to \( n \). Box \( k \) contains \( k \) balls numbered from 1 to \( k \). We randomly choose a box and then a ball from that box. Let \( X \) be the box number and \( Y \) be the ball number.
\begin{enumerate}
    \item Determine the joint distribution of \( (X, Y) \).
    \item Calculate \( P(X = Y) \).
    \item Determine the distribution of \( Y \) and \( E(Y) \).
\end{enumerate}

\begin{correctionbox}

\end{correctionbox}

% ----------- EXERCISE 6 -------------
\section{}
Let \( (X_i)_{i \in \mathbb{N}} \) be a sequence of independent Bernoulli random variables with parameter \( p \). Let \( Y_i = X_i X_{i+1} \).
\begin{enumerate}
    \item What is the distribution of \( Y_i \)?
    \item Let \( S_n = \sum_{i=1}^n Y_i \). Calculate \( E(S_n) \) and \( V(S_n) \).
\end{enumerate}

\begin{correctionbox}
    
\end{correctionbox}

% ----------- EXERCISE 7 -------------
\section{}
Let \( (X, Y) \) be a pair of discrete random variables with values in \( \mathbb{N}^2 \), with joint distribution:
\[
P([X = i] \cap [Y = j]) = \frac{a}{2^{i+1}j!}
\]
\begin{enumerate}
    \item Determine all possible values of \( a \).
    \item Determine the marginal distributions.
    \item Are \( X \) and \( Y \) independent?
\end{enumerate}

\begin{correctionbox}

\end{correctionbox}

% ----------- EXERCISE 8 -------------
\section{}
Let \( X_1, \dots, X_n \) be finite discrete random variables. Show that:
\[
E(X_1 + \dots + X_n) = \sum_{i=1}^n E(X_i).
\]

\begin{correctionbox}

\end{correctionbox}

\end{document}